%Revised version
\documentclass[12pt]{amsart}
\usepackage{latexsym,enumerate}
\usepackage{amssymb, xcolor,mathrsfs}
\usepackage{amsmath,amsthm,amsfonts,amssymb,latexsym, mathabx}
%\usepackage{amsfonts}
%\usepackage{showlabels}
\usepackage[pagebackref,colorlinks]{hyperref}
%\usepackage[T1]{fontenc}
%\usepackage[utf8]{inputenc}
%\usepackage{mathtools}
\usepackage{arydshln}

\headheight=7pt \textheight=574pt \textwidth=432pt \topmargin=14pt
\oddsidemargin=18pt \evensidemargin=18pt

\newtheorem{theorem}{Theorem}[section]
\newtheorem{lemma}[theorem]{Lemma}
\newtheorem{proposition}[theorem]{Proposition}
\newtheorem{corollary}[theorem]{Corollary}

\newtheorem{theorema}{Theorem}
\renewcommand{\thetheorema}{\Alph{theorema}}
\newtheorem{conj}[theorema]{Conjecture}

\newtheorem{consequence}[theorem]{Consequence}
\newtheorem{conjecture}[theorem]{Conjecture}
\newtheorem{problem}[theorem]{Problem}
\newtheorem{question}[theorem]{Question}
\newtheorem{MainTheorem}{Theorem}
\renewcommand{\theMainTheorem}{\Alph{MainTheorem}}

\theoremstyle{definition}
\newtheorem{definition}[theorem]{Definition}
\newtheorem{example}[theorem]{Example}
\newtheorem{exercise}[theorem]{Exercise}
\newtheorem{notation}[theorem]{Notation}
\newtheorem{remark}[theorem]{Remark}
\newtheorem{observation}[theorem]{Observation}
\newtheorem{hypothesis}[theorem]{Hypothesis}


\numberwithin{equation}{section}

\renewcommand{\theequation}{\arabic{section}.\arabic{equation}}
\newcommand{\GL}{{\mathrm {GL}}}
\newcommand{\PGL}{{\mathrm {PGL}}}
\newcommand{\SL}{{\mathrm {SL}}}
\newcommand{\PSL}{{\mathrm {PSL}}}
\newcommand{\POmega}{{\mathrm {P\Omega}}}


\newcommand{\GU}{{\mathrm {GU}}}
\newcommand{\PGU}{{\mathrm {PGU}}}
\newcommand{\SU}{{\mathrm {SU}}}
\newcommand{\PSU}{{\mathrm {PSU}}}
\newcommand{\U}{{\mathrm {U}}}

\newcommand{\GO}{{\mathrm {GO}}}
\newcommand{\Spin}{{\mathrm {Spin}}}
\newcommand{\SO}{{\mathrm {SO}}}
\newcommand{\PCO}{{\mathrm {PCO}}}
\newcommand{\CB}{\mathbf{C}}


\newcommand{\Sp}{{\mathrm {Sp}}}
\newcommand{\PSp}{{\mathrm {PSp}}}
\newcommand{\PCSp}{{\mathrm {PCSp}}}

\newcommand{\rk}{{\mathrm {rk}}}
\newcommand{\Ker}{\operatorname{Ker}}
\newcommand{\Aut}{{\mathrm {Aut}}}
\newcommand{\Out}{{\mathrm {Out}}}
\newcommand{\Mult}{{\mathrm {Mult}}}
%\newcommand{\Stab}{{\mathrm {Stab}}}
\newcommand{\Schur}{{\mathrm {Schur}}}
\newcommand{\Inn}{{\mathrm {Inn}}}
\newcommand{\Irr}{{\mathrm {Irr}}}
\newcommand{\IBR}{{\mathrm {IBr}}}
\newcommand{\IBRL}{{\mathrm {IBr}}_{\ell}}
\newcommand{\ord}{{\mathrm {ord}}}
\def\id{\mathop{\mathrm{ id}}\nolimits}
\renewcommand{\Im}{{\mathrm {Im}}}
\newcommand{\Ind}{{\mathrm {Ind}}}
\newcommand{\diag}{{\mathrm {diag}}}
\newcommand{\soc}{{\mathrm {soc}}}
\newcommand{\Soc}{{\mathrm {Soc}}}
\newcommand{\End}{{\mathrm {End}}}
\newcommand{\sol}{{\mathrm {sol}}}
\newcommand{\Hom}{{\mathrm {Hom}}}
\newcommand{\Mor}{{\mathrm {Mor}}}
\def\rank{\mathop{\mathrm{ rank}}\nolimits}
\newcommand{\Syl}{{\mathrm {Syl}}}
\newcommand{\St}{{\mathsf {St}}}
\newcommand{\Stab}{{\mathrm {Stab}}}
\newcommand{\Tr}{{\mathrm {Tr}}}
\newcommand{\tr}{{\mathrm {tr}}}
\newcommand{\cl}{{\mathrm {cl}}}
\newcommand{\Cl}{{\mathrm {Cl}}}
\newcommand{\cd}{{\mathrm {cd}}}
\newcommand{\lcm}{{\mathrm {lcm}}}
\newcommand{\Gal}{{\rm Gal}}
\newcommand{\Spec}{{\mathrm {Spec}}}
\newcommand{\ad}{{\mathrm {ad}}}
\newcommand{\Sym}{{\mathrm {Sym}}}
\newcommand{\Alt}{{\mathrm {Alt}}}
\newcommand{\Char}{{\mathrm {Char}}}
\newcommand{\Lin}{{\mathrm {Lin}}}
\newcommand{\sign}{{\mathrm {sgn}}}

\newcommand{\pr}{{\mathrm {pr}}}
\newcommand{\rad}{{\mathrm {rad}}}
\newcommand{\abel}{{\mathrm {abel}}}
\newcommand{\codim}{{\mathrm {codim}}}
\newcommand{\ind}{{\mathrm {ind}}}
\newcommand{\Res}{{\mathrm {Res}}}
\newcommand{\Ann}{{\mathrm {Ann}}}
\newcommand{\CC}{{\mathbb C}}
\newcommand{\RR}{{\mathbb R}}
\newcommand{\QQ}{{\mathbb Q}}
\newcommand{\ZZ}{{\mathbb Z}}
\newcommand{\NN}{{\mathbb N}}
\newcommand{\EE}{{\mathbb E}}
\newcommand{\PP}{{\mathbb P}}
\newcommand{\GG}{{\mathbb G}}
\newcommand{\HH}{{\mathbb HH}}

\newcommand{\SSS}{{\mathbb S}}
\newcommand{\AAA}{{\mathbb A}}
\newcommand{\HA}{\hat{{\mathbb A}}}
\newcommand{\FF}{{\mathbb F}}
\newcommand{\FP}{\mathbb{F}_{p}}
\newcommand{\FQ}{\mathbb{F}_{q}}
\newcommand{\FQB}{\mathbb{F}_{q}^{\bullet}}
\newcommand{\FQA}{\mathbb{F}_{q^{2}}}
\newcommand{\KK}{{\mathbb K}}
\newcommand{\EC}{\mathcal{E}}
\newcommand{\XC}{\mathcal{G}}
\newcommand{\HC}{\mathcal{H}}
\newcommand{\AC}{\mathcal{A}}
\newcommand{\OC}{\mathcal{O}}
\newcommand{\PC}{\mathcal{P}}

\newcommand{\OL}{O^{\ell'}}
\newcommand{\IL}{\mathcal{I}}
\newcommand{\CL}{\mathcal{C}}
\newcommand{\RL}{\mathcal{R}}
\newcommand{\TL}{\mathcal{T}}
\newcommand{\TC}{\mathcal{T}}
\newcommand{\GC}{\mathcal{G}}
%\newcommand{\GCR}{\mathcal{R}}
\newcommand{\eps}{\epsilon}
\newcommand{\LC}{\mathcal{L}}
\newcommand{\M}{\mathcal{M}}
%\newcommand{\HC}{\mathcal{H}}
%\newcommand{\bL}{\bar{L}}
\newcommand{\Om}{\Omega}
\newcommand{\tn}{\hspace{0.5mm}^{t}\hspace*{-1mm}}
\newcommand{\ta}{\hspace{0.5mm}^{2}\hspace*{-0.2mm}}
\newcommand{\tb}{\hspace{0.5mm}^{3}\hspace*{-0.2mm}}
\newcommand{\p}{^{\prime}}
\newcommand{\Center}{\mathbf{Z}}
\newcommand{\Centralizer}{\mathbf{C}}
\newcommand{\Normalizer}{\mathbf{N}}
\newcommand{\acd}{\mathsf{acd}}
\newcommand{\acde}{\mathsf{acd}_{\mathrm {even}}}
\newcommand{\acdo}{\mathsf{acd}_{\mathrm {odd}}}
\newcommand{\odd}{\mathrm{odd}}

\newcommand{\wt}{\widetilde}

\newcommand{\height}{\mathbf{ht}}
\newcommand{\lev}{\mathrm{\mathbf{lev}}}
\newcommand{\Bl}{\mathrm{\mathrm{Bl}}}


\newcommand{\acs}{\mathrm{acs}}
\newcommand{\bC}{{\mathbf{C}}}
\newcommand{\bG}{{\mathbf{G}}}
\newcommand{\bS}{{\mathbf{S}}}
\newcommand{\bH}{{\mathbf{H}}}
\newcommand{\bT}{{\mathbf{T}}}
\newcommand{\bL}{{\mathbf{L}}}
\newcommand{\bO}{{\mathbf{O}}}
\newcommand{\bN}{{\mathbf{N}}}
\newcommand{\bB}{{\mathbf{B}}}
\newcommand{\bU}{{\mathbf{U}}}
\newcommand{\NB}{{\mathbf{N}}}
\newcommand{\bZ}{{\mathbf{Z}}}
\newcommand{\bX}{{\mathbf{X}}}
\newcommand{\Al}{\textup{\textsf{A}}}
\newcommand{\Sy}{\textup{\textsf{S}}}
\newcommand{\tw}[1]{{}^#1}
\renewcommand{\mod}{\bmod \,}

\newcommand{\bff}{\mathbf{f}}

\newcommand{\pcore}{\mathbf{O}}
\newcommand{\GammaL}{\operatorname{\Gamma L}}
\newcommand{\IBr}{\operatorname{IBr}}
\newcommand{\bga}{\boldsymbol{\gamma}}
\newcommand{\blam}{\boldsymbol{\lambda}}
\newcommand{\lam}{\lambda}
\makeatletter \@namedef{subjclassname@2020}{\textup{2020}
Mathematics Subject Classification} \makeatother

\def\nor{\trianglelefteq\,}

\newcommand{\hungcomment}{\textcolor{red}}


\title[Hall $\pi$-subgroups and characters of $\pi'$-degree]{Hall $\pi$%%
%% Annotation 1, 1/30 on page 1 [ ]
%% Caret: "UBGROUPS AND<Caret> CHARACTERS "
%% Comment: ""
%%
%⭣ ⭣ ⭣
-subgroups and characters %%
%⭡ ⭡ ⭡ END of annotation 1
of $\pi'$-degree}

\author[Eugenio Giannelli]{Eugenio Giannelli}
\address{Dipartimento di Matematica e Informatica U. Dini, Firenze,
Italy}
\email{eugenio.giannelli@unifi.it}

\author[Nguyen N. Hung]{Nguyen N. Hung}
\address{Department of Mathematics, The University of Akron, Akron,
OH 44325, USA}
\email{hungnguyen@uakron.edu}


\thanks{The first author's research is funded by: the European Union Next
Generation EU, M4C1, CUP B53D23009410006, PRIN 2022 - 2022PSTWLB
Group Theory and Applications; and the INDAM-GNSAGA Project CUP
E53C24001950001. The second author gratefully acknowledges support
from the AMS--Simons Research Enhancement Grant (AWD-000167 AMS). We
thank Gunter Malle for several helpful comments on an earlier
version, particularly regarding
Proposition~\ref{prop:linear-unitary}}



\subjclass[2020]{Primary 20C15, 20C30, 20C33, 20D20}
%
\keywords{Characters, fields of values, Hall $\pi$-subgroups,
characters of $\pi'$-degree.}

%\date{}

\begin{document}
\begin{abstract}
%%
%% Annotation 2, 2/30 on page 1 [ ]
%% Caret: "ct. We study<Caret> the relatio"
%% Comment: ""
%%
%⭣ ⭣ ⭣
We study the relationship between the existence of Hall
$\pi$-subgroups and that of irreducible characters %%
%⭡ ⭡ ⭡ END of annotation 2
%%
%% Annotation 3, 3/30 on page 1 [ ]
%% Caret: "work extends<Caret> a result of"
%% Comment: ""
%%
%% Annotation 4, 4/30 on page 1 [ ]
%% Caret: "cters of π1-<Caret>degree with "
%% Comment: ""
%%
%% Annotation 5, 5/30 on page 1 [ ]
%% Caret: "racters of π<Caret>1-degree wit"
%% Comment: ""
%%
%% Annotation 6, 6/30 on page 1 [ ]
%% Caret: "egree with p<Caret>rescribed fi"
%% Comment: ""
%%
%⭣ ⭣ ⭣
of $\pi'$-degree
with prescribed fields of values in finite groups. This work extends
a result of Navarro and Tiep %%
%⭡ ⭡ ⭡ END of annotations 3, 4, 5, 6
from a single odd prime to multiple odd
primes.
\end{abstract}

\maketitle

%\tableofcontents

%%%%%%%%%%%%%%%%%%%%%%%%
%%%%%%%%%%%%%%%%%%%%%%%

\section{Introduction}

A well-known conjecture of R.~Gow, proved by G.~Navarro and
P.~H.~%%
%% Annotation 7, 7/30 on page 1 [ ]
%% Caret: "duction  A w<Caret>ell-known co"
%% Comment: ""
%%
%% Annotation 8, 8/30 on page 1 [ ]
%% Caret: "tion  A well<Caret>-known conje"
%% Comment: ""
%%
%⭣ ⭣ ⭣
Tiep~\cite{Navarro-Tiep2}, asserts that if %%
%⭡ ⭡ ⭡ END of annotations 7, 8
$G$ is %%
%% Annotation 9, 9/30 on page 1 [ ]
%% Caret: "-known conje<Caret>cture of R. "
%% Comment: ""
%%
%⭣ ⭣ ⭣
a finite
group %%
%⭡ ⭡ ⭡ END of annotation 9
of even order, then $G$ possesses a nontrivial (complex)
irreducible character of odd degree whose values are rational. This
result was subsequently generalized to all primes by the same
authors, as follows.

\begin{theorem}[\cite{Navarro-Tiep1}, Theorem~A]\label{thm:Navarro-Tiep}
Let $G$ be a finite group of order divisible by a prime $p$%%
%% Annotation 10, 10/30 on page 1 [ ]
%% Caret: "  Theorem 1.<Caret>1 ([NT06], T"
%% Comment: ""
%%
%⭣ ⭣ ⭣
. Then
$G$ %%
%⭡ ⭡ ⭡ END of annotation 10
%%
%% Annotation 11, 11/30 on page 1 [ ]
%% Caret: "1.1 ([NT06],<Caret> Theorem A)."
%% Comment: ""
%%
%⭣ ⭣ ⭣
has a nontrivial %%
%⭡ ⭡ ⭡ END of annotation 11
irreducible character of degree not divisible
by $p$ whose values lie %%
%% Annotation 12, 12/30 on page 1 [ ]
%% Caret: "in Qpe2πi{pq<Caret>.  In [GHSV2"
%% Comment: ""
%%
%% Annotation 13, 13/30 on page 1 [ ]
%% Caret: "ues lie in Q<Caret>pe2πi{pq.  I"
%% Comment: ""
%%
%⭣ ⭣ ⭣
in $\QQ(e^{2\pi i/p})$. %%
%⭡ ⭡ ⭡ END of annotations 12, 13
\end{theorem}

In~\cite{GHSV21}, together with A.\,A.~Schaeffer Fry and C.~Vallejo,
we attempted to extend this theorem from a single prime $p$ to a set
$\pi$ of two primes. Among other results, it was shown that %%
%% Annotation 14, 14/30 on page 1 [ ]
%% Caret: "t if π `` t2,<Caret> pu and gcdp"
%% Comment: ""
%%
%% Annotation 15, 15/30 on page 1 [ ]
%% Caret: "f π `` t2, pu<Caret> and gcdp|G|"
%% Comment: ""
%%
%% Annotation 16, 16/30 on page 1 [ ]
%% Caret: "|G|, 2pq {\k a} 1<Caret>, then G pos"
%% Comment: ""
%%
%% Annotation 17, 17/30 on page 1 [ ]
%% Caret: "G|, 2pq {\k a} 1,<Caret> then G poss"
%% Comment: ""
%%
%% Annotation 18, 18/30 on page 1 [ ]
%% Caret: " {\k a} 1, then G<Caret> possesses a"
%% Comment: ""
%%
%% Annotation 19, 19/30 on page 1 [ ]
%% Caret: "pq {\k a} 1, then<Caret> G possesses"
%% Comment: ""
%%
%% Annotation 20, 20/30 on page 1 [ ]
%% Caret: "in Qpe2πi{pq<Caret>. (Here, a c"
%% Comment: ""
%%
%% Annotation 21, 21/30 on page 1 [ ]
%% Caret: "n Qpe2πi{pq.<Caret> (Here, a ch"
%% Comment: ""
%%
%% Annotation 22, 22/30 on page 1 [ ]
%% Caret: "ntained in Q<Caret>pe2πi{pq. (H"
%% Comment: ""
%%
%⭣ ⭣ ⭣
if
$\pi=\{2,p\}$ and $\gcd(|G|,2p)>1$, then $G$ possesses a nontrivial
irreducible character of $\pi'$-degree whose values are contained in
$\QQ(e^{2\pi i/p})$. (Here, a character $\chi$ is said to have
$\pi'$-degree if $\chi(1)$ is not divisible by any prime in %%
%⭡ ⭡ ⭡ END of annotations 14, 15, 16, 17, 18, 19, 20, 21, 22
$\pi$.)
Unfortunately, this phenomenon does not hold in general when $\pi$
consists of two odd primes. For example, as pointed out in
\cite[Proposition~4.1]{GHSV21}, the Tits group ${}^2F_4(2)'$ has no
nontrivial irreducible character of %%
%% Annotation 27, 27/30 on page 1 [ ]
%% Caret: " of t3, 5u1-<Caret>degree with "
%% Comment: ""
%%
%% Annotation 28, 28/30 on page 1 [ ]
%% Caret: ", 5u1-degree<Caret> with values"
%% Comment: ""
%%
%⭣ ⭣ ⭣
$\{3,5\}'$-degree with %%
%⭡ ⭡ ⭡ END of annotations 27, 28
values in
$\QQ(e^{2\pi i/15})$.

In this paper, we propose a different extension -- perhaps a more
natural one -- of Theorem~\ref{thm:Navarro-Tiep} that works for
an arbitrary set of odd primes. In the following, we write $\pi(G)$
for the set of prime divisors of $|G|$. Recall that a subgroup $H\le
G$ is called %%
%% Annotation 29, 29/30 on page 1 [ ]
%% Caret: "alled a Hall<Caret> π-subgroup "
%% Comment: ""
%%
%% Annotation 30, 30/30 on page 1 [ ]
%% Caret: "all π-subgro<Caret>up of G if π"
%% Comment: ""
%%
%⭣ ⭣ ⭣
a Hall $\pi$-subgroup of %%
%⭡ ⭡ ⭡ END of annotations 29, 30
$G$ %%
%% Annotation 23, 23/30 on page 1 [ ]
%% Caret: "p of G if πp<Caret>Hq {\v D} π and |"
%% Comment: ""
%%
%% Annotation 24, 24/30 on page 1 [ ]
%% Caret: " of G if πpH<Caret>q {\v D} π and |G"
%% Comment: ""
%%
%% Annotation 25, 25/30 on page 1 [ ]
%% Caret: "q {\v D} π and |G<Caret> : H| is not"
%% Comment: ""
%%
%% Annotation 26, 26/30 on page 1 [ ]
%% Caret: "{\v D} π and |G :<Caret> H| is not d"
%% Comment: ""
%%
%⭣ ⭣ ⭣
if $\pi(H)\subseteq \pi$
and $|G:H|$ is %%
%⭡ ⭡ ⭡ END of annotations 23, 24, 25, 26
not divisible by any prime in $\pi$.

\begin{theorema}\label{thm:main}
Let $\pi$ be a set of odd primes and $G$ a finite group possessing a
nontrivial Hall $\pi$-subgroup. Then $G$ has a nontrivial
$\pi'$-degree irreducible character with values in $\QQ(e^{2\pi
i/p})$ for some $p\in\pi$.
\end{theorema}



Theorem~\ref{thm:main} may be viewed as a generalization of
Theorem~\ref{thm:Navarro-Tiep}. In fact, when $p$ is odd,
Theorem~\ref{thm:Navarro-Tiep} follows from the case $\pi=\{p\}$ of
Theorem~\ref{thm:main} together with the first Sylow theorem.

If we assume that $2\in \pi$, then the conclusion of
Theorem~\ref{thm:main} remains true when %%
%% Annotation 31, 1/10 on page 2 [ ]
%% Caret: "hen |π| {d'} 2,<Caret> as noted ab"
%% Comment: ""
%%
%⭣ ⭣ ⭣
$|\pi|\leq 2$, as %%
%⭡ ⭡ ⭡ END of annotation 31
noted
above, but fails in general once $|\pi|\geq 3$. For example, if $G$
is any non-abelian simple group and %%
%% Annotation 32, 2/10 on page 2 [ ]
%% Caret: "nd π `` πpGq,<Caret> then the tr"
%% Comment: ""
%%
%⭣ ⭣ ⭣
$\pi=\pi(G)$, then %%
%⭡ ⭡ ⭡ END of annotation 32
the trivial
character $1_G$ is the only irreducible character of $G$ with
$\pi'$-degree. There are also counterexamples with $\pi\varsubsetneq
\pi(G)$, such as $(G,\pi)=(\Al_7,\{2,3,5\})$. At present, we do not
have a conceptual explanation for why the prime $2$ is special in
this context.

Similar to Gow's conjecture and the results in
\cite{Navarro-Tiep1,GHSV21}, Theorem~\ref{thm:main} admits a clean
reduction to finite simple groups. Accordingly, the main body of
this paper is devoted to proving the result for such groups. This
requires some new work on the relationship between the existence of
Hall $\pi$-subgroups and $\pi'$-degree irreducible characters in
simple groups of Lie type.


%%%%%%%%%%%%%%%%%%%%%%%
%%%%%%%%%%%%%%%%%%%%%%%
%\section{Alternating groups}
%For notational simplicity, we write $\QQ(\zeta_k)$ for the $k$-th
%cyclotomic field $\QQ(e^{2\pi i/k})$. As usual, $\Irr(G)$ denotes
%the set of all irreducible characters of $G$, and $\Irr_{\pi'}(G)$
%denotes the subset consisting of characters whose degrees are not
%divisible by any prime in $\pi$. Given $x,y\in\mathbb{Z}$ we let $[x,y]:=\{z\in\mathbb{Z}\ |\ x\leq z\leq y\}$.
%
%\begin{lemma}\label{lem:noHall}
%Let $n\in\mathbb{N}_{\geq 5}$ and let $\pi\subseteq [1,n]$ be a set of prime numbers. Assume that $2\notin\pi$.
%Then the alternating group $\Al_n$ admits a Hall $\pi$-subgroup if and only if $|\pi|=1$.
%\end{lemma}
%\begin{proof}
%If $|\pi|=1$ the statement is obviously implied by Sylow theory.
%Let us now assume that $\Al_n$ admits a Hall $\pi$-subgroup $H$ and let us suppose for a contradiction that $|\pi|\geq 2$.
%In particular let $\pi=\{p_1,p_2,\ldots, p_t\}$ for some odd primes $p_1<p_2<\cdots<p_t$.
%We proceed by distinguishing two cases, depending on the solvability of $H$.
%
%\smallskip
%
%%\noindent\textit{Case} (i).
%\noindent $\bullet$ Let us first assume that $H$ is a solvable group. By \cite{Hall1} we know that $H$ admits a $\{p_1,p_2\}$-Hall subgroup $K$. It follows that $K$ is a solvable $\{p_1,p_2\}$-Hall subgroup of $\Al_n$. Moreover, since $2\notin \{p_1,p_2\}$ we also have that $K$ is a solvable $\{p_1,p_2\}$-Hall subgroup of $\Sy_n$. Using \cite[Theorem A4]{Hall2}, we deduce that $|K|$ is even and therefore that $2\in \{p_1,p_2\}\subseteq \pi$. This clearly contradicts our hypothesis.
%
%
%
%\smallskip
%
%\noindent $\bullet$ Let us now consider the case where $H$ is a non-solvable Hall $\pi$-subgroup of $\Al_n$. Again, since $2\notin \pi$ we have that $H$ is a non-solvable Hall $\pi$-subgroup of $\Sy_n$.
%This family of subgroups of symmetric groups is completely described in \cite{Thompson}. In particular, we deduce that either $H=\Sy_n$ or $n$ is a prime number and $H=\Sy_{n-1}$. In both cases we would have that $2\in\pi$, contradicting our hypothesis.
%\end{proof}
%
%The above Lemma \ref{lem:noHall} immediately implies that the following statement is vacuously true.
%
%\begin{theorem}\label{thm:alternating}
%Let $n\in\ZZ_{\ge 5}$ and $\pi$ be a set of odd primes such that
%$\Al_n$ has a Hall $\pi$-subgroup. Then there exists
%$1_{\Al_n}\neq\chi\in\Irr_{\pi'}(\Al_n)$ with
%$\QQ(\chi)\subseteq\QQ(\zeta_p)$ for some $p\in\pi$.
%\end{theorem}
%
%Notice that the $(2\notin \pi)$-hypothesis in the statement of Theorem \ref{thm:alternating} can not be removed.
%In fact, as we already pointed out in the introduction, $\Al_7$ posses a Hall $\{2,3,5\}$-subgroup but the only $\{2,3,5\}'$-degree irreducible character of $\Al_7$ is the trivial character.
%Actually, this observation can be extended to any prime larger than $7$.
%Let $p\geq 11$ be a fixed prime number, and let $\pi$ be the set consisting of all prime numbers strictly smaller than $p$.
%In this case $\Al_{p-1}$ is a Hall $\pi$-subgroup of $\Al_p$ but $\mathrm{Irr}_{\pi'}(\Al_p)=\{1_{\Al_p}\}$. In fact any non-trivial irreducible character $\chi$ of $\pi'$-degree of $\Al_p$ would satisfy $\chi(1)=p$. This is not possible because for any $\zeta\in\mathrm{Irr}(\Sy_p)$ we have that either $\zeta(1)\leq p-1$ or that $\zeta(1)\geq \frac{p(p-3)}{2}$, by \cite{Rasala}. Since $p-3> 4$ we have that for any $\theta\in\mathrm{Irr}(\Al_p)$ we have that either $\theta(1)\leq p-1$ or that $\theta(1)>p$.


%%%%%%%%%%%%%%%%%%%%%%%
%%%%%%%%%%%%%%%%%%%%%%%
\section{Non-Abelian Simple Groups}

This section is devoted to proving Theorem~\ref{thm:main} for all
non-abelian simple groups. We first fix some notation. For a
positive integer $k$, we %%
%% Annotation 37, 7/10 on page 2 [ ]
%% Caret: "ζk :`` e2πi{k<Caret> and we writ"
%% Comment: ""
%%
%⭣ ⭣ ⭣
let $\zeta_k:=e^{2\pi i/k}$ and %%
%⭡ ⭡ ⭡ END of annotation 37
we write $\QQ(\zeta_k)$ for the $k$th
cyclotomic field. As usual, $\Irr(G)$ denotes
the set of all irreducible characters of a group $G$, and
$\Irr_{\pi'}(G)$ denotes the subset consisting of those characters
whose degrees are not divisible by any prime in $\pi$. For
$\chi\in\Irr(G)$, we write $\QQ(\chi)$ for the field of values of
$\chi$, that is, the smallest extension of $\QQ$ containing all
values of $\chi$. Finally, for integers $x\le y$, we let
$[x,y]:=\{\,z\in\ZZ \mid x\le z\le y\,\}$.

As mentioned above, we aim at proving the following.

\begin{theorem}\label{thm:simple}
Let $S$ be a non-abelian simple group and $\pi$ be a set of odd
primes such that $S$ has a Hall $\pi$-subgroup. Then there exists
%%
%% Annotation 36, 6/10 on page 2 [ ]
%% Caret: "χ P Irrπ1pSq<Caret> such that Q"
%% Comment: ""
%%
%⭣ ⭣ ⭣
$1_S\neq\chi\in\Irr_{\pi'}(S)$ such %%
%⭡ ⭡ ⭡ END of annotation 36
%%
%% Annotation 33, 3/10 on page 2 [ ]
%% Caret: "Qpχq {\v D} Qpζpq<Caret> for some p "
%% Comment: ""
%%
%⭣ ⭣ ⭣
that
$\QQ(\chi)\subseteq\QQ(\zeta_p)$ for %%
%⭡ ⭡ ⭡ END of annotation 33
some $p\in\pi$.
\end{theorem}


Hall $\pi$-subgroups for a set $\pi$ of odd primes are relatively
common in simple groups of Lie type, but are much more restricted in
alternating and sporadic groups. We begin by treating the
alternating and sporadic cases.

\begin{lemma}\label{lem:noHall}
%%
%% Annotation 35, 5/10 on page 2 [ ]
%% Caret: " Let n P N{\v e}5<Caret> and let π {\v D}"
%% Comment: ""
%%
%⭣ ⭣ ⭣
Let $n\in\mathbb{N}_{\geq 5}$ and %%
%⭡ ⭡ ⭡ END of annotation 35
let $\pi\subseteq [1,n]$ be a set
of prime numbers. Assume that $2\notin\pi$. Then the alternating
group $\Al_n$ admits a Hall $\pi$-subgroup if and only if $|\pi|=1$.
\end{lemma}

\begin{proof}
%%
%% Annotation 34, 4/10 on page 2 [ ]
%% Caret: ". If |π| `` 1<Caret> the stateme"
%% Comment: ""
%%
%⭣ ⭣ ⭣
If $|\pi|=1$ the %%
%⭡ ⭡ ⭡ END of annotation 34
statement is obviously implied by the Sylow theory.
For the other implication, assume that $\Al_n$ admits a Hall
$\pi$-subgroup $H$ and suppose for a contradiction that $|\pi|\geq
2$. %%
%% Annotation 38, 8/10 on page 2 [ ]
%% Caret: " . . . , ptu<Caret> for some od"
%% Comment: ""
%%
%% Annotation 39, 9/10 on page 2 [ ]
%% Caret: "rticular let<Caret> π `` tp1, p2"
%% Comment: ""
%%
%% Annotation 40, 10/10 on page 2 [ ]
%% Caret: "n particular<Caret> let π `` tp1"
%% Comment: ""
%%
%⭣ ⭣ ⭣
In particular let $\pi=\{p_1,p_2,\ldots, p_t\}$ for %%
%⭡ ⭡ ⭡ END of annotations 38, 39, 40
some odd
primes $p_1<p_2<\cdots<p_t$. %We proceed by distinguishing two cases,
%depending on the solvability of $H$.
Notice that $H$ is solvable, by the Feit--Thompson theorem
\cite{Thompson63}. By \cite{Hall1} we then know that $H$ admits %%
%% Annotation 41, 1/30 on page 3 [ ]
%% Caret: " a tp1, p2u-<Caret>Hall subgrou"
%% Comment: ""
%%
%% Annotation 42, 2/30 on page 3 [ ]
%% Caret: "s a tp1, p2u<Caret>-Hall subgro"
%% Comment: ""
%%
%% Annotation 43, 3/30 on page 3 [ ]
%% Caret: "p1, p2u-Hall<Caret> subgroup K."
%% Comment: ""
%%
%⭣ ⭣ ⭣
a
$\{p_1,p_2\}$-Hall subgroup %%
%⭡ ⭡ ⭡ END of annotations 41, 42, 43
$K$. It follows that $K$ is a solvable
$\{p_1,p_2\}$-Hall subgroup of $\Al_n$. Moreover, since $2\notin
\{p_1,p_2\}$ we also have that $K$ is a solvable $\{p_1,p_2\}$-Hall
subgroup of the symmetric group $\Sy_n$. Using \cite[Theorem
A4]{Hall2}, we deduce that $|K|$ is even, and %%
%% Annotation 44, 4/30 on page 3 [ ]
%% Caret: "2 P tp1, p2u<Caret> {\v D} π. This c"
%% Comment: ""
%%
%% Annotation 45, 5/30 on page 3 [ ]
%% Caret: " therefore 2<Caret> P tp1, p2u "
%% Comment: ""
%%
%⭣ ⭣ ⭣
therefore $2\in
\{p_1,p_2\}\subseteq \pi$. %%
%⭡ ⭡ ⭡ END of annotations 44, 45
This clearly contradicts our hypothesis.
%
%
%%\noindent\textit{Case} (i).
%(i) First assume that $H$ is a solvable group. By \cite{Hall1} we
%know that $H$ admits a $\{p_1,p_2\}$-Hall subgroup $K$. It follows
%that $K$ is a solvable $\{p_1,p_2\}$-Hall subgroup of $\Al_n$.
%Moreover, since $2\notin \{p_1,p_2\}$ we also have that $K$ is a
%solvable $\{p_1,p_2\}$-Hall subgroup of the symmetric group $\Sy_n$. Using \cite[Theorem
%A4]{Hall2}, we deduce that $|K|$ is even, and therefore $2\in
%\{p_1,p_2\}\subseteq \pi$. This clearly contradicts our hypothesis.
\end{proof}




We note that the use of the Feit--Thompson theorem can be avoided by
instead relying on the known classification of non-solvable Hall
subgroups of symmetric groups. Suppose that $H$ is a non-solvable
Hall $\pi$-subgroup of $\Al_n$. Since $2 \notin \pi$, it follows
that $H$ is also a %%
%% Annotation 46, 6/30 on page 3 [ ]
%% Caret: "olvable Hall<Caret> π-subgroup "
%% Comment: ""
%%
%⭣ ⭣ ⭣
non-solvable Hall $\pi$%%
%⭡ ⭡ ⭡ END of annotation 46
-subgroup of $\Sy_n$. This
family of subgroups of symmetric groups is completely described in
\cite{Thompson}. In particular, either $H = \Sy_n$ or $n$ is a prime
number and $H = \Sy_{n-1}$. In both cases, we would have $2 \in
\pi$, which again contradicts our hypothesis.


%The above Lemma \ref{lem:noHall} immediately implies that the following statement is vacuously true.

\begin{proposition}\label{prop:alternatingandsporadic}
Theorem~\ref{thm:simple} holds when $S$ is an alternating group, a
sporadic group, or the %%
%% Annotation 47, 7/30 on page 3 [ ]
%% Caret: "oup 2F4p2q1.<Caret>  Proof The "
%% Comment: ""
%%
%% Annotation 48, 8/30 on page 3 [ ]
%% Caret: "e Tits group<Caret> 2F4p2q1.  P"
%% Comment: ""
%%
%⭣ ⭣ ⭣
Tits group ${}^2F_4(2)'$.
%Let $n\in\ZZ_{\ge 5}$ and $\pi$ be a set of odd primes such that
%$\Al_n$ has a Hall $\pi$-subgroup. Then there exists
%$1_{\Al_n}\neq\chi\in\Irr_{\pi'}(\Al_n)$ with
%$\QQ(\chi)\subseteq\QQ(\zeta_p)$ for some $p\in\pi$.
\end{proposition}

\begin{proof}
The %%
%⭡ ⭡ ⭡ END of annotations 47, 48
case of alternating groups follows from Lemma~\ref{lem:noHall}
and Theorem~\ref{thm:Navarro-Tiep}. We now consider the sporadic
groups and the Tits group. By \cite[Theorem~6.14]{G86}, if $S$ has a
Hall $\pi$-subgroup for a set $\pi$ of odd primes, %%
%% Annotation 70, 30/30 on page 3 [ ]
%% Caret: ", then |π| {d'}<Caret> 2. The resu"
%% Comment: ""
%%
%⭣ ⭣ ⭣
then $|\pi| \le
2$. %%
%⭡ ⭡ ⭡ END of annotation 70
The result then follows from \cite[Theorem~2.1]{GHSV21}, except
possibly in the cases $(S,\pi)=({}^2F_4(2)',\{3,5\})$, $(J_4,
\{23,43\})$, or $(J_4,\{29,43\})$. However, a direct check using
\cite{Atl1} shows that %%
%% Annotation 67, 27/30 on page 3 [ ]
%% Caret: "that in each<Caret> of these ca"
%% Comment: ""
%%
%⭣ ⭣ ⭣
in each of %%
%⭡ ⭡ ⭡ END of annotation 67
these cases, $S$ does not have a
Hall %%
%% Annotation 66, 26/30 on page 3 [ ]
%% Caret: "-subgroup. □<Caret>  R k 2 4 Th"
%% Comment: ""
%%
%% Annotation 68, 28/30 on page 3 [ ]
%% Caret: "l π-subgroup<Caret>. □  Remark "
%% Comment: ""
%%
%% Annotation 69, 29/30 on page 3 [ ]
%% Caret: " π-subgroup.<Caret> □  Remark 2"
%% Comment: ""
%%
%⭣ ⭣ ⭣
$\pi$-subgroup.
\end{proof}

%In particular we have that Theorem \ref{thm:simple} holds whenever $S$ is a simple alternating group.

\begin{remark}
The %%
%⭡ ⭡ ⭡ END of annotations 66, 68, 69
hypothesis $2\notin \pi$ in
Proposition~\ref{prop:alternatingandsporadic} can not be removed.
Indeed, as already pointed out in the introduction, $\Al_7$
possesses a Hall $\{2,3,5\}$-subgroup %%
%% Annotation 56, 16/30 on page 3 [ ]
%% Caret: "bgroup but t<Caret>he only t2, "
%% Comment: ""
%%
%⭣ ⭣ ⭣
but the only %%
%⭡ ⭡ ⭡ END of annotation 56
$\{2,3,5\}'$-degree irreducible character of $\Al_7$ is the trivial
%%
%% Annotation 54, 14/30 on page 3 [ ]
%% Caret: "racter. This<Caret> observation"
%% Comment: ""
%%
%% Annotation 57, 17/30 on page 3 [ ]
%% Caret: "cter. This o<Caret>bservation c"
%% Comment: ""
%%
%⭣ ⭣ ⭣
character. This observation can %%
%⭡ ⭡ ⭡ END of annotations 54, 57
be extended to any prime larger than
$7$. %%
%% Annotation 55, 15/30 on page 3 [ ]
%% Caret: ". Let p {\v e} 11<Caret> be a fixed "
%% Comment: ""
%%
%⭣ ⭣ ⭣
Let $p\geq 11$ be %%
%⭡ ⭡ ⭡ END of annotation 55
a fixed %%
%% Annotation 58, 18/30 on page 3 [ ]
%% Caret: "ixed prime n<Caret>umber, and l"
%% Comment: ""
%%
%⭣ ⭣ ⭣
prime number, and %%
%⭡ ⭡ ⭡ END of annotation 58
let $\pi$ be the
set consisting of all prime numbers strictly %%
%% Annotation 53, 13/30 on page 3 [ ]
%% Caret: "smaller than<Caret> p. In this "
%% Comment: ""
%%
%% Annotation 59, 19/30 on page 3 [ ]
%% Caret: "aller than p<Caret>. In this ca"
%% Comment: ""
%%
%⭣ ⭣ ⭣
smaller than $p$. In %%
%⭡ ⭡ ⭡ END of annotations 53, 59
this case $\Al_{p-1}$ is a Hall $\pi$-subgroup of $\Al_p$ but
$\mathrm{Irr}_{\pi'}(\Al_p)=\{1_{\Al_p}\}$. In fact any non-trivial
irreducible character $\chi$ of $\pi'$-degree of $\Al_p$ would
satisfy $\chi(1)=p$. %%
%% Annotation 52, 12/30 on page 3 [ ]
%% Caret: "`` p. This is<Caret> not possibl"
%% Comment: ""
%%
%⭣ ⭣ ⭣
This is not %%
%⭡ ⭡ ⭡ END of annotation 52
possible because for any
$\zeta\in\mathrm{Irr}(\Sy_p)$ we have either $\zeta(1)\leq p-1$ %%
%% Annotation 50, 10/30 on page 3 [ ]
%% Caret: " ppp´3q  2 ,<Caret> by [Ras77]."
%% Comment: ""
%%
%% Annotation 51, 11/30 on page 3 [ ]
%% Caret: "p1q {\v e} ppp´3q<Caret>  2 , by [Ra"
%% Comment: ""
%%
%⭣ ⭣ ⭣
or
$\zeta(1)\geq \frac{p(p-3)}{2}$, by %%
%⭡ ⭡ ⭡ END of annotations 50, 51
\cite{Rasala}. Since $p-3> 4$ we
have that for any $\theta\in\mathrm{Irr}(\Al_p)$, either
$\theta(1)\leq p-1$ or $\theta(1)>p$%%
%% Annotation 60, 20/30 on page 3 [ ]
%% Caret: "<Caret>e now turn t"
%% Comment: ""
%%
%⭣ ⭣ ⭣
.
\end{remark}

\smallskip

We now %%
%⭡ ⭡ ⭡ END of annotation 60
turn to the proof of Theorem~\ref{thm:simple} for simple
groups of Lie type $S \neq {}^2F_4(2)'$. By this, we mean simple
groups of the form $S = G/\bZ(G)$%%
%% Annotation 61, 21/30 on page 3 [ ]
%% Caret: "{ZpGq, where<Caret> G :`` GF is "
%% Comment: ""
%%
%% Annotation 62, 22/30 on page 3 [ ]
%% Caret: "here G :`` GF<Caret> is the grou"
%% Comment: ""
%%
%% Annotation 63, 23/30 on page 3 [ ]
%% Caret: ":`` GF is the<Caret> group of fi"
%% Comment: ""
%%
%⭣ ⭣ ⭣
, where $G := \bG^F$ is the group %%
%⭡ ⭡ ⭡ END of annotations 61, 62, 63
of fixed points of a simple, simply connected algebraic group $\bG$
%%
%% Annotation 64, 24/30 on page 3 [ ]
%% Caret: "defined over<Caret> an algebrai"
%% Comment: ""
%%
%⭣ ⭣ ⭣
defined over an %%
%⭡ ⭡ ⭡ END of annotation 64
algebraically closed field of characteristic $\ell$,
under a Steinberg %%
%% Annotation 65, 25/30 on page 3 [ ]
%% Caret: "domorphism F<Caret> of G.  "
%% Comment: ""
%%
%⭣ ⭣ ⭣
endomorphism $F$ of %%
%⭡ ⭡ ⭡ END of annotation 65
$\bG$.
%; see
%\cite[Theorem~24.17]{malletesterman}.

The %%
%% Annotation 49, 9/30 on page 3 [ ]
%% Caret: "he case ℓR π<Caret> is easy.  "
%% Comment: ""
%%
%⭣ ⭣ ⭣
case $\ell\notin\pi$ is %%
%⭡ ⭡ ⭡ END of annotation 49
easy.

\begin{proposition}\label{prop:ell-notin-pi}
Theorem \ref{thm:simple} holds when $S$ is a simple group of Lie
type in characteristic $\ell$ and $\ell\notin\pi$.
\end{proposition}

\begin{proof}
The Steinberg character $\St_G$ of $G$ is trivial on $\bZ(G)$ and
thus can be viewed as a character of $S$. Moreover, $\St_G$ is
rational-valued and has degree a power of $\ell$; see
\cite[Proposition~3.4.10]{GeckMalle}.
\end{proof}

We therefore focus on the case $\ell\in\pi$. Let $\bB$ be an
$F$-stable Borel subgroup of $\bG$ and $\bT$ a (maximally-split)
$F$-stable maximal torus of $\bG$ inside $\bB$. Let $\Phi$ be the
root system of $\bG$ with respect to $\bT$ and $\bB$ and $\Phi^+$
the set of corresponding positive roots. Let $\bU$ be the product of
the root subgroups corresponding to the roots in $\Phi^+$. This
$\bU$ is indeed the unipotent radical of $\bB$. We have
$U:=\bU^F\in\Syl_\ell(G)$ and $B=UT$, %%
%% Annotation 77, 7/9 on page 4 [ ]
%% Caret: "and T :`` TF;<Caret> see [Car85,"
%% Comment: ""
%%
%% Annotation 78, 8/9 on page 4 [ ]
%% Caret: " and T :`` TF<Caret>; see [Car85"
%% Comment: ""
%%
%% Annotation 79, 9/9 on page 4 [ ]
%% Caret: "here U :`` UF<Caret> and T :`` TF"
%% Comment: ""
%%
%⭣ ⭣ ⭣
where $U:=\bU^F$ and
$T:=\bT^F$; see %%
%⭡ ⭡ ⭡ END of annotations 77, 78, 79
\cite[\S1.9 and \S2.9]{Carter85}.

Let $(\bG^*,F^*)$ be the pair dual to $(\bG,F)$, and %%
%% Annotation 75, 5/9 on page 4 [ ]
%% Caret: "G˚ :`` G˚F ˚.<Caret> Note that S"
%% Comment: ""
%%
%% Annotation 76, 6/9 on page 4 [ ]
%% Caret: " G˚ :`` G˚F ˚<Caret>. Note that "
%% Comment: ""
%%
%⭣ ⭣ ⭣
set
$G^*:={\bG^*}^{F^*}$. Note %%
%⭡ ⭡ ⭡ END of annotations 75, 76
that $S=[G^*,G^*]$, as $S$ is simple. Let
$\bT^*$ be an $F^*$-stable maximal torus of $\bG^*$ that is dual to
$\bT$ in the sense of \cite[Proposition~4.3.1]{Carter85}. Let
$\bB^*$ be an %%
%% Annotation 74, 4/9 on page 4 [ ]
%% Caret: "˚ be an F ˚-<Caret>stable Borel"
%% Comment: ""
%%
%⭣ ⭣ ⭣
$F^*$-stable Borel %%
%⭡ ⭡ ⭡ END of annotation 74
subgroup of $\bG^*$ %%
%% Annotation 72, 2/9 on page 4 [ ]
%% Caret: "ontaining T˚<Caret>. Write T ˚ "
%% Comment: ""
%%
%⭣ ⭣ ⭣
containing
$\bT^*$. Write %%
%⭡ ⭡ ⭡ END of annotation 72
$T^*:={\bT^*}^{F^*}$ %%
%% Annotation 73, 3/9 on page 4 [ ]
%% Caret: " B˚ :`` B˚F ˚<Caret>. If U˚ deno"
%% Comment: ""
%%
%⭣ ⭣ ⭣
and $B^*:={\bB^*}^{F^*}$. If %%
%⭡ ⭡ ⭡ END of annotation 73
$\bU^*$ denotes the unipotent radical %%
%% Annotation 71, 1/9 on page 4 [ ]
%% Caret: "adical of B˚<Caret> and U ˚ :`` "
%% Comment: ""
%%
%⭣ ⭣ ⭣
of $\bB^*$ and %%
%⭡ ⭡ ⭡ END of annotation 71
$U^*:={\bU^*}^{F^*}$, then $B^*=U^*T^*$.


\begin{lemma}\label{lem:Gross}
Let $S$ be a simple group of Lie type in characteristic $\ell$.
Suppose that $\pi$ is a set of odd primes containing $\ell$ (in
particular, $\ell$ is odd) and that $S$ has a Hall $\pi$-subgroup.
Then $B^*$ contains a Hall $\pi$-subgroup of $G^*$.
\end{lemma}

\begin{proof}
Note that $\bZ(G)\le T\le B$. By \cite[Theorem~3.2]{G86}, the
quotient $B/\bZ(G)$, which may be regarded as a Borel subgroup of
$S$, contains a Hall $\pi$-subgroup of $S$. Equivalently, $|G:B| =
|S : B/\bZ(G)|$ is a $\pi'$-number. By Corollary~4.4.2 and
Proposition~4.4.4 of \cite{Carter85}, we have $|G^*|=|G|$ and
$|T^*|=|T|$, and it follows that $|U^*|=|U|$ and $|B|=|B^*|$. (Note
that $U$ and $U^*$ have orders equal to the $\ell$-parts of $|G|$
and $|G^*|$, respectively.) Now we have that $|G^*:B^*|$ is a
$\pi'$-number. As $B^*$ is solvable, it therefore contains a Hall
$\pi$-subgroup of $G^*$.
\end{proof}

\begin{lemma}\label{lem:Gross2}
Let $S\neq {}^2F_4(2)'$ be a simple group of Lie type in
characteristic $\ell$. To prove Theorem~\ref{thm:simple} for $S$, it
is sufficient to assume that $\pi\cap\pi(T^*)\neq\emptyset$.
\end{lemma}

\begin{proof}
By Proposition~\ref{prop:ell-notin-pi}, we may assume that
$\ell\in\pi$. If $\pi\cap \pi(S)=\{\ell\}$, then
Theorem~\ref{thm:simple} follows from
Theorem~\ref{thm:Navarro-Tiep}. We therefore may assume that
$|\pi\cap \pi(S)|\geq 2$. Since $S$ has a Hall $\pi$-subgroup, by
Lemma~\ref{lem:Gross}, the Borel subgroup $B^*$ contains a Hall
$\pi$-subgroup of $G^*$. In particular, $|\pi\cap \pi(B^*)|\geq 2$,
and hence there exists at least one prime lying in both $\pi$ and
$\pi(T^*)$.
\end{proof}



Recall that $S=G/\bZ(G)$, so the irreducible characters of $S$ are
precisely those irreducible characters of $G$ whose kernel contains
$\bZ(G)$. The set $\Irr(G)$ admits a natural partition into Lusztig
series $\EC(G,s)$ indexed by $G^\ast$-conjugacy classes of
semisimple elements $s\in G^\ast$. The series $\EC(G,s)$ consists of
those irreducible characters of $G$ that are constituents of some
Deligne--Lusztig character $R_{\bS}^{\bG}(\theta)$, where $\bS$ is
an $F$-stable maximal torus of $\bG$ and $\theta\in\Irr(\bS^F)$ is
such that the geometric conjugacy class of the pair $(\bS,\theta)$
corresponds to the $G^\ast$-conjugacy class containing $s$. See
\cite[\S2.6]{GeckMalle}.

We will look for the desired character among the so-called
\emph{semisimple characters}. Recall that \(\rho \in \Irr(G)\) is
called a semisimple character if the average value of \(\rho\) on
\(\mathfrak{C}^F\) is nonzero, where \(\mathfrak{C}\) is the
conjugacy class of \(\bG\) consisting of regular unipotent elements.
For each semisimple element \(s \in G^*\), the Lusztig series
\(\EC(G,s)\) contains one or more semisimple characters, all of
which have %%
%% Annotation 80, 1/5 on page 5 [ ]
%% Caret: ": CG˚psq|ℓ1.<Caret>  In fact he"
%% Comment: ""
%%
%⭣ ⭣ ⭣
degree
\[
|G^* : \bC_{G^*}(s)|_{\ell'}.
\]
In %%
%⭡ ⭡ ⭡ END of annotation 80
fact, when \(\bC_{\bG^*}(s)\) is connected, \(\EC(G,s)\) contains
exactly one semisimple character; see \cite[p.~171]{GeckMalle}. We
record a well-known fact about these characters below. Here we use
$\ord(g)$ to denote the order of a group element $g$.

\begin{lemma}\label{lem:1}
Assume the above notation. Assume furthermore that $\bC_{\bG^*}(s)$
is connected. Then the unique semisimple character, say $\chi_s$, in
the series \(\EC(G,s)\) is trivial on $\bZ(G)$ if and only if $s\in
S$. In such situation, $\QQ(\chi_s)\subseteq \QQ(\zeta_{\ord(s)})$.
\end{lemma}

\begin{proof}
The first part follows from, for instance,
\cite[Proposition~24.21]{malletesterman}, \cite[Lemma~2.2]{M07}, and
\cite[Lemma~5.8]{Hung22}. The latter part is
\cite[Lemma~4.2]{GHSV21}.
\end{proof}

We handle the linear and unitary groups separately. To unify
notation, we write $S=\PSL_n^\epsilon(q)$, where the superscript
$\epsilon=+1$ corresponds to the linear groups and $\epsilon=-1$ to
the unitary groups. We use analogous notation for the related
groups, for example $G=\SL_n^\epsilon(q)$ and
$G^*=\PGL_n^\epsilon(q)$. It is more convenient to first study the
characters of $\widetilde{G}:=\GL_n^\epsilon(q)$ and then analyze
those of $G=\SL_n^\epsilon(q)$ as irreducible constituents of
restricted characters. Note that $\widetilde{G}$ is self-dual, and
we will identify $\widetilde{G}$ with its dual group. Let
$\pi:\widetilde{G}\to G^*$ be the natural projection from
$\widetilde{G}$ to $G^*$.

Let $\widetilde{s}$ be a semisimple element of $\widetilde{G}$.
Since the ambient algebraic group of $\widetilde{G}$ has connected
center, the Lusztig series
$\mathcal{E}(\widetilde{G},\widetilde{s})$ contains a unique
semisimple character, which we denote by $\chi_{\widetilde{s}}$.

\begin{lemma}\label{lem:numberofconstituents}
The number of irreducible constituents of the restriction of
semisimple character $\chi_{\widetilde{s}}$ to $G$ divides
$\gcd(\ord(\widetilde{s}),n,q-\epsilon)$.
\end{lemma}

\begin{proof}
Set $s:=\pi(\widetilde{s})$. Recall that
$\chi_{\widetilde{s}}(1)=|\widetilde{G}:\bC_{\widetilde{G}}(\widetilde{s})|_{\ell'}$.
The restriction of $\chi_{\widetilde{s}}$ from $\widetilde{G}$ to
$G$ is multiplicity-free and its irreducible constituents are
precisely the semisimple characters of the Lusztig series
\(\EC(G,s)\) by \cite[Corollary~2.6.18]{GeckMalle}. Each of these
constituent has degree $|G^* : \bC_{G^*}(s)|_{\ell'}$. Therefore,
the number of irreducible constituents of the restriction %%
%% Annotation 81, 2/5 on page 5 [ ]
%% Caret: "<Caret> |π´1pCG˚psq"
%% Comment: ""
%%
%% Annotation 82, 3/5 on page 5 [ ]
%% Caret: "C rGprsq|ℓ1.<Caret>  | rG : C r"
%% Comment: ""
%%
%% Annotation 83, 4/5 on page 5 [ ]
%% Caret: "´1pCG˚psqq :<Caret> C rGprsq|ℓ1"
%% Comment: ""
%%
%% Annotation 84, 5/5 on page 5 [ ]
%% Caret: " : CG˚psq|ℓ1<Caret> `` |G : C rG"
%% Comment: ""
%%
%⭣ ⭣ ⭣
is
\[
\frac{|\widetilde{G}:\bC_{\widetilde{G}}(\widetilde{s})|_{\ell'}}{|G^*
:
\bC_{G^*}(s)|_{\ell'}}=\frac{|\widetilde{G}:\bC_{\widetilde{G}}(\widetilde{s})|_{\ell'}}{|\widetilde{G}
:
\pi^{-1}(\bC_{G^*}(s))|_{\ell'}}=|\pi^{-1}(\bC_{G^*}(s)):\bC_{\widetilde{G}}(\widetilde{s})|_{\ell'}.
\]




Let %%
%⭡ ⭡ ⭡ END of annotations 81, 82, 83, 84
$Z$ be the (cyclic) subgroup of the multiplicative group
$\FF_{q^2-1}^\times$ of order $q-\epsilon$. Consider the mapping
%%
%% Annotation 95, 11/13 on page 6 [ ]
%% Caret: "psqq Ñ Z def<Caret>ined by grsg"
%% Comment: ""
%%
%% Annotation 96, 12/13 on page 6 [ ]
%% Caret: "psqq Ñ Z def<Caret>ined by grsg"
%% Comment: ""
%%
%% Annotation 97, 13/13 on page 6 [ ]
%% Caret: "˚psqq Ñ Z de<Caret>fined by grs"
%% Comment: ""
%%
%⭣ ⭣ ⭣
$\kappa: \pi^{-1}(\bC_{G^*}(s))\to Z$ defined by %%
%⭡ ⭡ ⭡ END of annotations 95, 96, 97
$g\widetilde{s}g^{-1}=\kappa(g)\widetilde{s}$. This is indeed a
homomorphism with $\Ker(\kappa)=\bC_{\widetilde{G}}(\widetilde{s})$.
It follows that the index
$|\pi^{-1}(\bC_{G^*}(s)):\bC_{\widetilde{G}}(\widetilde{s})|$ is
equal to the order of the image $\Im(\kappa)$ of $\kappa$.
Therefore, the lemma follows once we show %%
%% Annotation 85, 1/13 on page 6 [ ]
%% Caret: ", n, q ´ ϵq.<Caret>  First, it "
%% Comment: ""
%%
%⭣ ⭣ ⭣
that
\[
|\Im(\kappa)| \text{ divides }
\gcd(\ord(\widetilde{s}),n,q-\epsilon).
\]


First, %%
%⭡ ⭡ ⭡ END of annotation 85
it %%
%% Annotation 86, 2/13 on page 6 [ ]
%% Caret: " it is clear<Caret> that |Impκq"
%% Comment: ""
%%
%% Annotation 87, 3/13 on page 6 [ ]
%% Caret: "at |Impκq| |<Caret> pq ´ ϵq. No"
%% Comment: ""
%%
%% Annotation 88, 4/13 on page 6 [ ]
%% Caret: "q| | pq ´ ϵq<Caret>. Now let g "
%% Comment: ""
%%
%⭣ ⭣ ⭣
is clear that $|\Im(\kappa)|\mid (q-\epsilon)$. Now %%
%⭡ ⭡ ⭡ END of annotations 86, 87, 88
let
$g$ be an arbitrary element of $\pi^{-1}(\bC_{G^*}(s))$. We have
$\det(\widetilde{s})=\det(g\widetilde{s}g^{-1})=\det(\kappa(g)\widetilde{s})=\kappa(g)^n\det(\widetilde{s})$,
and hence $\kappa(g)^n=1$; equivalently, $\ord(\kappa(g))\mid n$.
Finally, we observe that
$\ord(\widetilde{s})=\ord(g\widetilde{s}g^{-1})=\ord(\kappa(g)\widetilde{s})=\lcm(\ord(\kappa(g)),\widetilde{s})$,
which implies that $\ord(\kappa(g))\mid \ord(\widetilde{s})$.
%\[\{\delta^{(q-1)/p}, \delta^{(q-1)(p-1)/p},
%1^{\,n-2}\}=\{\kappa(g)\delta^{(q-1)/p},
%\kappa(g)\delta^{(q-1)(p-1)/p}, \kappa(g)^{\,n-2}\},\] where
%$\kappa(g)\in\bZ(\widetilde{G})$ is identified with the
%corresponding element in $\FF_q^\times$. This happens if and only if
%$\kappa(g)=1$. We have shown that $\pi^{-1}(\bC_{G^*}(s))=
%\bC_{\widetilde{G}}(\widetilde{s})$, and the claim follows.
%Now we have
%\[
%|G^*:\bC_{G^*}(s)|_{\ell'}=|G^*:\bC_{\widetilde{G}}(\widetilde{s})/\bZ(\widetilde{G})|_{\ell'}
%=|\widetilde{G}:\bC_{\widetilde{G}}(\widetilde{s})|_{\ell'},
%\]
%and therefore the semisimple characters in the series
%$\mathcal{E}(G,s)$ have the same degree as the (only) semisimple
%character $\chi_{\widetilde{s}}$ in
%$\mathcal{E}(\widetilde{G},\widetilde{s})$. However, by
%\cite[Corollary~2.6.18]{GeckMalle}, semisimple characters in
%$\mathcal{E}(G,s)$ are irreducible constituents of the restriction
%of $\chi_{\widetilde{s}}$ from $\widetilde{G}$ to $G$. We deduce
%that $\mathcal{E}(G,s)$ has a unique semisimple character, namely
%the restriction of $\chi_{\widetilde{s}}$ to $G$. Let us denote this
%restriction by $\psi$.
\end{proof}

The proof of Lemma~\ref{lem:numberofconstituents} also yeilds the
following well-known result; however, we only require the %%
%% Annotation 89, 5/13 on page 6 [ ]
%% Caret: "t statement.<Caret>  Le a 2 10 "
%% Comment: ""
%%
%% Annotation 90, 6/13 on page 6 [ ]
%% Caret: "st statement<Caret>.  Le a 2 10"
%% Comment: ""
%%
%⭣ ⭣ ⭣
first
statement.

\begin{lemma}\label{lem:2}
Let %%
%⭡ ⭡ ⭡ END of annotations 89, 90
$s := \pi(\widetilde{s}) \in G^*$, where $\widetilde{s}$ is a
semisimple element of $\widetilde{G}$. Suppose that the multiset of
eigenvalues of $\widetilde{s}$ is not invariant under multiplication
by every nontrivial root of unity. Then the centralizer
$\bC_{\bG^*}(s)$ is connected. Moreover, if the multiset of
eigenvalues of $\widetilde{s}$ is not invariant under multiplication
by every nontrivial element of the subgroup of $\FF_{q^2-1}^\times$
of order $q-\epsilon$, then $\chi_{\widetilde{s}}$ restricts
irreducibly to $G$.
\end{lemma}

\begin{proof}
For every $g \in \pi^{-1}\bigl(\bC_{G^*}(s)\bigr)$, we have
$g\widetilde{s}g^{-1} = \kappa(g)\widetilde{s}$, and hence
$g\widetilde{s}g^{-1}$ and $\widetilde{s}$ have the same
eigenvalues. Therefore, if the multiset of eigenvalues of
$\widetilde{s}$ is not invariant under multiplication by every
nontrivial element of the subgroup of $\FF_{q^2-1}^\times$ of order
$q-\epsilon$, the homomorphism $\kappa$ must have trivial image,
which proves the second statement.

A similar argument, now in the setting of algebraic groups
$\widetilde{\bG} := \GL^\epsilon_n(\overline{\FF}_{\ell})$ and
${\bG^*} := \PGL^\epsilon_n(\overline{\FF}_{\ell})$ instead, shows
that if the multiset of eigenvalues of $\widetilde{s}$ is not
invariant under multiplication by every nontrivial root of unity,
then $\bC_{\bG^*}(s)$ is the image of the connected group
$\bC_{\widetilde{\bG}}(\widetilde{s})$ under $\pi$ and hence is
itself connected.
\end{proof}


\begin{proposition}\label{prop:linear-unitary}
Theorem~\ref{thm:simple} holds for the simple groups $S=\PSL_n(q)$
and $\PSU_n(q)$, where $n\geq 2$ and $q$ is a prime power.
\end{proposition}

\begin{proof}
Recall that $T^*$ is a maximal torus contained in a Borel subgroup
$B^*$ of $G^*$. By Lemma~\ref{lem:Gross2}, we may assume that $\pi$
and $\pi(T^*)$ share a common prime, say $p$. In the linear case, we
%%
%% Annotation 91, 7/13 on page 6 [ ]
%% Caret: "`` pq ´ 1qn´1<Caret>,  so p divi"
%% Comment: ""
%%
%% Annotation 92, 8/13 on page 6 [ ]
%% Caret: " pq ´ 1qn´1,<Caret>  so p divid"
%% Comment: ""
%%
%⭣ ⭣ ⭣
have
\[
|T^*| = (q-1)^{n-1},
\]
so %%
%⭡ ⭡ ⭡ END of annotations 91, 92
$p$ divides $q-1$. In the unitary %%
%% Annotation 93, 9/13 on page 6 [ ]
%% Caret: " if n is eve<Caret>n,  "
%% Comment: ""
%%
%% Annotation 94, 10/13 on page 6 [ ]
%% Caret: " ´ 1qpn´2q{2<Caret>pq ´ 1q, if "
%% Comment: ""
%%
%⭣ ⭣ ⭣
case,
\[
|T^*| =
\begin{cases}
(q^2-1)^{(n-1)/2}, & \text{if $n$ is odd},\\[2pt]
(q^2-1)^{(n-2)/2}(q-1), & \text{if $n$ is even},
\end{cases}
\]
and %%
%⭡ ⭡ ⭡ END of annotations 93, 94
hence $p$ divides $q^2-1$ (see \cite[p.~74]{Carter85}). Since
$B^*$ contains a Hall $\pi$-subgroup of $G^*$ by
Lemma~\ref{lem:Gross}, a comparison of the orders of $G^*$ and $B^*$
shows that $p$ cannot divide $q+1$. Therefore, as in the linear
case, we again conclude that $p$ divides $q-1$.



We first consider the case $\epsilon=1$. Let $\delta$ be a generator
of $\FF_{q}^\times$, and define the semisimple %%
%% Annotation 98, 1/50 on page 7 [ ]
%% Caret: " diag  P rG.<Caret>  Let s deno"
%% Comment: ""
%%
%⭣ ⭣ ⭣
element
\[
\widetilde{s} := \diag\!\bigl( \delta^{(q-1)/p},
\delta^{(q-1)(p-1)/p}, 1^{\,n-2} \bigr) \in \widetilde{G}.
\]
Let %%
%⭡ ⭡ ⭡ END of annotation 98
$s$ denote the image $\pi(\widetilde{s})$ of $\widetilde{s}$
%%
%% Annotation 99, 2/50 on page 7 [ ]
%% Caret: " π : rG Ñ G˚<Caret>. Then  ordp"
%% Comment: ""
%%
%% Annotation 100, 3/50 on page 7 [ ]
%% Caret: "G Ñ G˚. Then<Caret>  ordpsq `` o"
%% Comment: ""
%%
%% Annotation 101, 4/50 on page 7 [ ]
%% Caret: "ordprsq `` p.<Caret>  Taking con"
%% Comment: ""
%%
%⭣ ⭣ ⭣
under $\pi:\widetilde{G}\to G^*$. Then
\[
\ord(s) = \ord(\widetilde{s}) = p.
\]
Taking %%
%⭡ ⭡ ⭡ END of annotations 99, 100, 101
conjugation if necessary, we may assume that $s\in T^*$.
Moreover, as $\det(\widetilde{s})=\delta^{q-1}=1$%%
%% Annotation 102, 5/50 on page 7 [ ]
%% Caret: "have  rs P G<Caret> and s P S X"
%% Comment: ""
%%
%% Annotation 103, 6/50 on page 7 [ ]
%% Caret: "s P S X T ˚.<Caret> Note that t"
%% Comment: ""
%%
%% Annotation 104, 7/50 on page 7 [ ]
%% Caret: " s P S X T ˚<Caret>. Note that "
%% Comment: ""
%%
%⭣ ⭣ ⭣
, we have
\[
\widetilde{s}\in G \quad\text{and}\quad s\in S\cap T^*.
\]

Note %%
%⭡ ⭡ ⭡ END of annotations 102, 103, 104
that the case $(n,p)=(3,3)$ cannot occur. (Otherwise, by
Lemma~\ref{lem:Gross}, the Borel subgroup $B^*$, of order
$q^3(q-1)^2$, would contain a Sylow $3$-subgroup of $G^*$ whose
order $q^3(q^2-1)(q^3-1)$ has larger $3$-part, leading to a
contradiction.) It is then easy to check that the multiset of
eigenvalues of $\widetilde{s}$ is not invariant under multiplication
by every nontrivial root of unity. By Lemma~\ref{lem:2}, %%
%% Annotation 142, 45/50 on page 7 [ ]
%% Caret: "2.10, it fol<Caret>lows that CG"
%% Comment: ""
%%
%% Annotation 143, 46/50 on page 7 [ ]
%% Caret: " 2.10, it fo<Caret>llows that C"
%% Comment: ""
%%
%% Annotation 144, 47/50 on page 7 [ ]
%% Caret: ".10, it foll<Caret>ows that CG˚"
%% Comment: ""
%%
%⭣ ⭣ ⭣
it follows
that %%
%⭡ ⭡ ⭡ END of annotations 142, 143, 144
$\bC_{\bG^*}(s)$ is connected, and we let $\chi_s$ be the
unique semisimple character in the series $\EC(G,s)$. Using
Lemma~\ref{lem:1}, we deduce %%
%% Annotation 105, 8/50 on page 7 [ ]
%% Caret: "sqq `` Qpζpq.<Caret>  Moreover s"
%% Comment: ""
%%
%⭣ ⭣ ⭣
that
\[\QQ(\chi_{s})\subseteq\QQ(\zeta_{\ord(s)})=\QQ(\zeta_p).\] Moreover, %%
%⭡ ⭡ ⭡ END of annotation 105
since $s\in T^*$ and $T^*$ is abelian, we %%
%% Annotation 106, 9/50 on page 7 [ ]
%% Caret: " ˚ {d'} CG˚psq.<Caret>  It follows"
%% Comment: ""
%%
%% Annotation 107, 10/50 on page 7 [ ]
%% Caret: " |G˚ : T ˚|,<Caret> and hence  "
%% Comment: ""
%%
%% Annotation 108, 11/50 on page 7 [ ]
%% Caret: "follows that<Caret>  |G˚ : CG˚p"
%% Comment: ""
%%
%% Annotation 109, 12/50 on page 7 [ ]
%% Caret: ".  It follow<Caret>s that  |G˚ "
%% Comment: ""
%%
%% Annotation 110, 13/50 on page 7 [ ]
%% Caret: "  It follows<Caret> that  |G˚ :"
%% Comment: ""
%%
%% Annotation 111, 14/50 on page 7 [ ]
%% Caret: "˚psq.  It fo<Caret>llows that  "
%% Comment: ""
%%
%% Annotation 112, 15/50 on page 7 [ ]
%% Caret: "psq.  It fol<Caret>lows that  |"
%% Comment: ""
%%
%% Annotation 113, 16/50 on page 7 [ ]
%% Caret: " T ˚|, and h<Caret>ence  χsp1q "
%% Comment: ""
%%
%% Annotation 114, 17/50 on page 7 [ ]
%% Caret: "T ˚|, and he<Caret>nce  χsp1q ``"
%% Comment: ""
%%
%% Annotation 115, 18/50 on page 7 [ ]
%% Caret: " ˚|, and hen<Caret>ce  χsp1q `` "
%% Comment: ""
%%
%% Annotation 116, 19/50 on page 7 [ ]
%% Caret: "˚|, and henc<Caret>e  χsp1q `` |"
%% Comment: ""
%%
%% Annotation 117, 20/50 on page 7 [ ]
%% Caret: "|, and hence<Caret>  χsp1q `` |G"
%% Comment: ""
%%
%% Annotation 118, 21/50 on page 7 [ ]
%% Caret: " : T ˚|, and<Caret> hence  χsp1"
%% Comment: ""
%%
%% Annotation 119, 22/50 on page 7 [ ]
%% Caret: "˚ : T ˚|, an<Caret>d hence  χsp"
%% Comment: ""
%%
%% Annotation 120, 23/50 on page 7 [ ]
%% Caret: "G˚ : T ˚|, a<Caret>nd hence  χs"
%% Comment: ""
%%
%⭣ ⭣ ⭣
have \[T^*\le
\bC_{\bG^*}(s).\] It follows that \[|G^*:\bC_{\bG^*}(s)| \text{
divides } |G^*:T^*|,\] and hence
\[
\chi_s(1)=|G^*:\bC_{G^*}(s)|_{\ell'} \text{ divides }
|G^*:T^*|_{\ell'}.
\]
As %%
%⭡ ⭡ ⭡ END of annotations 106, 107, 108, 109, 110, 111, 112, 113, 114, 115, 116, 117, 118, 119, 120
$|G^*:T^*|_{\ell'}=|G^*:B^*|$ is a $\pi'$-number %%
%% Annotation 141, 44/50 on page 7 [ ]
%% Caret: "umber by Lem<Caret>ma 2.6, we c"
%% Comment: ""
%%
%⭣ ⭣ ⭣
by
Lemma~\ref{lem:Gross}, we %%
%⭡ ⭡ ⭡ END of annotation 141
conclude that $\chi_s(1)$ is also a
$\pi'$-number. Also, from the fact that $s\in S=[G^*,G^*]$, we know
that $\chi_{s}$ is trivial %%
%% Annotation 121, 24/50 on page 7 [ ]
%% Caret: "trivial on Z<Caret>pGq. This co"
%% Comment: ""
%%
%⭣ ⭣ ⭣
on $\bZ({G})$. %%
%⭡ ⭡ ⭡ END of annotation 121
This completes the proof
for linear %%
%% Annotation 122, 25/50 on page 7 [ ]
%% Caret: "r groups.  I<Caret>t remains to"
%% Comment: ""
%%
%% Annotation 123, 26/50 on page 7 [ ]
%% Caret: "roups.  It r<Caret>emains to co"
%% Comment: ""
%%
%% Annotation 124, 27/50 on page 7 [ ]
%% Caret: "ups.  It rem<Caret>ains to cons"
%% Comment: ""
%%
%% Annotation 125, 28/50 on page 7 [ ]
%% Caret: "oups.  It re<Caret>mains to con"
%% Comment: ""
%%
%% Annotation 126, 29/50 on page 7 [ ]
%% Caret: "ps.  It rema<Caret>ins to consi"
%% Comment: ""
%%
%% Annotation 127, 30/50 on page 7 [ ]
%% Caret: ".  It remain<Caret>s to conside"
%% Comment: ""
%%
%% Annotation 128, 31/50 on page 7 [ ]
%% Caret: "It remains t<Caret>o consider t"
%% Comment: ""
%%
%% Annotation 129, 32/50 on page 7 [ ]
%% Caret: "emains to co<Caret>nsider the c"
%% Comment: ""
%%
%% Annotation 130, 33/50 on page 7 [ ]
%% Caret: "remains to c<Caret>onsider the "
%% Comment: ""
%%
%% Annotation 131, 34/50 on page 7 [ ]
%% Caret: "mains to con<Caret>sider the ca"
%% Comment: ""
%%
%% Annotation 132, 35/50 on page 7 [ ]
%% Caret: "ns to consid<Caret>er the case "
%% Comment: ""
%%
%% Annotation 133, 36/50 on page 7 [ ]
%% Caret: "ns to consid<Caret>er the case "
%% Comment: ""
%%
%% Annotation 134, 37/50 on page 7 [ ]
%% Caret: "ins to consi<Caret>der the case"
%% Comment: ""
%%
%% Annotation 135, 38/50 on page 7 [ ]
%% Caret: "ains to cons<Caret>ider the cas"
%% Comment: ""
%%
%% Annotation 136, 39/50 on page 7 [ ]
%% Caret: "s to conside<Caret>r the case ϵ"
%% Comment: ""
%%
%% Annotation 137, 40/50 on page 7 [ ]
%% Caret: "o consider t<Caret>he case ϵ `` "
%% Comment: ""
%%
%% Annotation 138, 41/50 on page 7 [ ]
%% Caret: "consider the<Caret> case ϵ `` ´1"
%% Comment: ""
%%
%% Annotation 139, 42/50 on page 7 [ ]
%% Caret: " consider th<Caret>e case ϵ `` ´"
%% Comment: ""
%%
%% Annotation 140, 43/50 on page 7 [ ]
%% Caret: "nsider the c<Caret>ase ϵ `` ´1 a"
%% Comment: ""
%%
%⭣ ⭣ ⭣
groups.


It remains to consider the case $\epsilon=-1$ %%
%⭡ ⭡ ⭡ END of annotations 122, 123, 124, 125, 126, 127, 128, 129, 130, 131, 132, 133, 134, 135, 136, 137, 138, 139, 140
and $p \mid (q-1)$.
Let \(\widetilde{t}\) be a semisimple element of order \(p\) in a
maximal torus of order $q^2-1$ of \(\GU_2(q)\). Since \(p \mid
(q-1)\) and \(p\) is odd, we have \(\gcd(p,q+1)=1\), and hence
\(\widetilde{t}\in \SU_2(q)\). Define the semisimple element
\[
\widetilde{s}:=\diag\left(\widetilde{t}, I_{n-2}\right)\in
\widetilde{G},
\]
and consider the corresponding semisimple character
\(\chi_{\widetilde{s}}\in \Irr(\widetilde{G})\). Note that
$\widetilde{s}$ belongs to a maximal torus of $\widetilde{G}$ of
order $(q^2-1)^{(n-1)/2}(q+1)$ if $n$ is odd, and of %%
%% Annotation 145, 48/50 on page 7 [ ]
%% Caret: "er pq2´1qn{2<Caret>  if n is ev"
%% Comment: ""
%%
%% Annotation 146, 49/50 on page 7 [ ]
%% Caret: "f n is even.<Caret>  "
%% Comment: ""
%%
%% Annotation 147, 50/50 on page 7 [ ]
%% Caret: "´1qn{2  if n<Caret> is even.  "
%% Comment: ""
%%
%⭣ ⭣ ⭣
order
$(q^2-1)^{n/2}$ if $n$ is even.

Note %%
%⭡ ⭡ ⭡ END of annotations 145, 146, 147
also that \(\ord(\widetilde{s})=\ord(\widetilde{t})=p\), which
does not divide $q+1$. Lemma~\ref{lem:numberofconstituents} then
implies that \(\chi_{\widetilde{s}}\) restricts irreducibly to
\(G\). Arguments similar to those in the linear case show that this
restriction is trivial on \(\bZ(G)\), has \(\pi'\)-degree, and takes
values in \(\QQ(\zeta_p)\). This completes the proof.
\end{proof}


We can now finish the proof of Theorem~\ref{thm:simple}.

\begin{proof}[Proof of Theorem~\ref{thm:simple}]
By Propositions~\ref{prop:alternatingandsporadic} and
\ref{prop:linear-unitary}, and the classification of finite simple
groups, we may assume %%
%% Annotation 161, 14/14 on page 8 [ ]
%% Caret: " S ‰ 2F4p2q1<Caret> is a simple"
%% Comment: ""
%%
%⭣ ⭣ ⭣
that $S \neq {}^2F_4(2)'$ is %%
%⭡ ⭡ ⭡ END of annotation 161
a simple group of
Lie type not of type $A$. By Lemma~\ref{lem:Gross2}, there exists at
least one prime lying in both $\pi$ and $\pi(T^*)$. As before, let
$p$ be such a prime.

Note that
\[
|T^*:(T^*\cap S)| = |T^*S:S| \text{ divides } |G^*:S|
\]
and $|G^*:S|$ is the order of the group of \emph{diagonal
automorphisms} of $S$. On the other hand, the order of $T^*$ is
given %%
%% Annotation 148, 1/14 on page 8 [ ]
%% Caret: " ``  q|o| ´ 1<Caret>  oPO  "
%% Comment: ""
%%
%⭣ ⭣ ⭣
by
\[
|T^*|=\prod_{\mathfrak{o}\in \mathcal{O}}
\bigl(q^{|\mathfrak{o}|}-1\bigr),
\]
where %%
%⭡ ⭡ ⭡ END of annotation 148
$\mathcal{O}$ is the set of $F^*$-orbits on the simple roots
of the root system of $\bG^*$ associated with $\bB^*$ and $\bT^*$,
and $q$ is the absolute value of the eigenvalues of $F^*$ on the
character group of $\bT^*$; see \cite[p.~74]{Carter85}. A
straightforward case-by-case check, using \cite[Table~5]{Atl1} for
the size of the group of diagonal automorphisms and
\cite[\S1.19]{Carter85} for the sizes of the $F^*$-orbits, reveals
that if an odd prime $p$ is a divisor of $|G^*:S|$, then indeed
$(|G^*:S|_p)^2$ divides $|T^*|$, and it follows that $p$ also
divides $|T^*\cap S|$. (The only exception is the case $S=\PSU_3(q)$
with $(q+1)_3=3$, but this was already excluded using
Proposition~\ref{prop:linear-unitary}.) Therefore, we may and will
assume from now on that \[p\in \pi\cap \pi(T^*\cap S).\]
Consequently, there exists a semisimple %%
%% Annotation 154, 7/14 on page 8 [ ]
%% Caret: "ement s P G˚<Caret> such that  "
%% Comment: ""
%%
%% Annotation 155, 8/14 on page 8 [ ]
%% Caret: "lement s P G<Caret>˚ such that "
%% Comment: ""
%%
%⭣ ⭣ ⭣
element $s\in G^*$ such %%
%⭡ ⭡ ⭡ END of annotations 154, 155
%%
%% Annotation 153, 6/14 on page 8 [ ]
%% Caret: " ordpsq `` p.<Caret>  Suppose fi"
%% Comment: ""
%%
%⭣ ⭣ ⭣
that
\[
s\in T^*\cap S \quad\text{and}\quad \ord(s)=p.
\]

Suppose %%
%⭡ ⭡ ⭡ END of annotation 153
first that $p\nmid |\bZ(\bG)|$. Then $\bC_{\bG^*}(s)$ is
connected by \cite[Corollary~4.6]{Borel70}. Arguing as in the proof
of Proposition~\ref{prop:linear-unitary} and applying
Lemmas~\ref{lem:Gross} and \ref{lem:1}, we have that the unique
semisimple character $\chi_s$ in $\EC(G,s)$ has $\pi'$-degree, is
trivial on $\bZ(G)$, and satisfies $\QQ(\chi_s)\subseteq
\QQ(\zeta_p)$.

Next suppose that $p\mid |\bZ(\bG)|$. By
\cite[Theorem~1.12.5]{GLS98} and the assumption that $p$ is odd, we
are left with only the %%
%% Annotation 152, 5/14 on page 8 [ ]
%% Caret: "e case p `` 3<Caret> and G to be"
%% Comment: ""
%%
%⭣ ⭣ ⭣
case $p=3$ and %%
%⭡ ⭡ ⭡ END of annotation 152
$\bG$ to be of type $E_6$.
Here, %%
%% Annotation 158, 11/14 on page 8 [ ]
%% Caret: "G `` Eϵ  6pqq<Caret> with ϵ P t˘"
%% Comment: ""
%%
%⭣ ⭣ ⭣
$G=E_6^\epsilon(q)$ with %%
%⭡ ⭡ ⭡ END of annotation 158
%%
%% Annotation 149, 2/14 on page 8 [ ]
%% Caret: "ith ϵ P t˘1u<Caret>, where ϵ `` "
%% Comment: ""
%%
%% Annotation 150, 3/14 on page 8 [ ]
%% Caret: "th ϵ P t˘1u,<Caret> where ϵ `` 1"
%% Comment: ""
%%
%⭣ ⭣ ⭣
$\epsilon\in\{\pm1\}$, where %%
%⭡ ⭡ ⭡ END of annotations 149, 150
%%
%% Annotation 151, 4/14 on page 8 [ ]
%% Caret: " where ϵ `` 1<Caret> corresponds"
%% Comment: ""
%%
%⭣ ⭣ ⭣
$\epsilon=1$ corresponds %%
%⭡ ⭡ ⭡ END of annotation 151
to the untwisted groups %%
%% Annotation 159, 12/14 on page 8 [ ]
%% Caret: "s and ϵ `` ´1<Caret> to the twis"
%% Comment: ""
%%
%% Annotation 160, 13/14 on page 8 [ ]
%% Caret: "ps and ϵ `` ´<Caret>1 to the twi"
%% Comment: ""
%%
%⭣ ⭣ ⭣
and $\epsilon=-1$
to %%
%⭡ ⭡ ⭡ END of annotations 159, 160
the twisted groups, and $|\bZ(G)|=\gcd(3,q-\epsilon)=3$.

As mentioned above, all the semisimple characters in the Lusztig
series $\EC(G,s)$ indexed by $s$ have the same degree
$|G^*:\bC_{G^*}(s)|_{\ell'}$. This is a $\pi'$-number due to the
fact that %%
%% Annotation 156, 9/14 on page 8 [ ]
%% Caret: "that s P T ˚<Caret>, as argued "
%% Comment: ""
%%
%% Annotation 157, 10/14 on page 8 [ ]
%% Caret: "hat s P T ˚,<Caret> as argued i"
%% Comment: ""
%%
%⭣ ⭣ ⭣
$s\in T^*$, as %%
%⭡ ⭡ ⭡ END of annotations 156, 157
argued in the proof of
Proposition~\ref{prop:linear-unitary}. Moreover, as noted in the
proof of \cite[Proposition~4.5]{GHSV21}, these semisimple characters
take integer values on unipotent elements and therefore have field
of values contained in $\QQ(\zeta_3)$, by \cite[Lemma~4.3]{GHSV21}.
Finally, since their degrees are coprime to $3$, their restrictions
%%
%% Annotation 181, 20/23 on page 9 [ ]
%% Caret: "ions to ZpGq<Caret> are multipl"
%% Comment: ""
%%
%⭣ ⭣ ⭣
to $\bZ(G)$ are %%
%⭡ ⭡ ⭡ END of annotation 181
multiples of the trivial character. (Let $\chi$ be
such a character and %%
%% Annotation 182, 21/23 on page 9 [ ]
%% Caret: "ZpGq `` χp1qα<Caret> for some α "
%% Comment: ""
%%
%⭣ ⭣ ⭣
let $\chi_{\bZ(G)}=\chi(1)\alpha$ for %%
%⭡ ⭡ ⭡ END of annotation 182
some
$\alpha\in\Irr(\bZ(G))$. %%
%% Annotation 183, 22/23 on page 9 [ ]
%% Caret: "1qαq `` αχp1q<Caret>, which impl"
%% Comment: ""
%%
%% Annotation 184, 23/23 on page 9 [ ]
%% Caret: "`` detpχZpGqq<Caret> `` detpχp1qα"
%% Comment: ""
%%
%⭣ ⭣ ⭣
Then
$\mathbf{1}_{\bZ(G)}=\det(\chi)_{\bZ(G)}=\det(\chi_{\bZ(G)})=\det(\chi(1)\alpha)=\alpha^{\chi(1)}$,
which %%
%⭡ ⭡ ⭡ END of annotations 183, 184
implies that the order of $\alpha$ divides
$(\chi(1),|\bZ(G)|)$.) This completes the proof.
\end{proof}

%%%%%%%%%%%%%%%%%%%%%%%
%%%%%%%%%%%%%%%%%%%%%%%
\section{A Proof of Theorem \ref{thm:main}}

We are now in the position to give a complete proof of our main result, as stated in the introduction.

A finite group admitting a Hall $\pi$-subgroup is often called an
\emph{$E_\pi$-group}. There is a large literature on the theory of
$E_\pi$-groups, including the results about simple groups we cited in
the previous section. We refer the reader to \cite{RV10} for the
latest results and relevant information. Note that the class of
$E_\pi$-groups is not closed under extensions and a subgroup of an
$E_\pi$-group might be no longer an $E_\pi$-group. Nevertheless, the
only fact we need is that every normal/subnormal subgroup and every
quotient of an $E_\pi$-group is an $E_\pi$-group. In fact, if $H$ is
a Hall $\pi$-subgroup of $G$ and $N\nor G$, then $H\cap N$ is a Hall
$\pi$-subgroup of $N$ %%
%% Annotation 178, 17/23 on page 9 [ ]
%% Caret: "f N and HN{N<Caret> is a Hall π"
%% Comment: ""
%%
%⭣ ⭣ ⭣
and $HN/N$ is %%
%⭡ ⭡ ⭡ END of annotation 178
%%
%% Annotation 179, 18/23 on page 9 [ ]
%% Caret: "{N is a Hall<Caret> π-subgroup "
%% Comment: ""
%%
%% Annotation 180, 19/23 on page 9 [ ]
%% Caret: " is a Hall π<Caret>-subgroup of"
%% Comment: ""
%%
%⭣ ⭣ ⭣
a Hall $\pi$-subgroup of %%
%⭡ ⭡ ⭡ END of annotations 179, 180
$G/N$.


Another fact we need is that %%
%% Annotation 176, 15/23 on page 9 [ ]
%% Caret: "hat if N {\IJ } G<Caret> such that |"
%% Comment: ""
%%
%⭣ ⭣ ⭣
if $N\nor G$ such %%
%⭡ ⭡ ⭡ END of annotation 176
that %%
%% Annotation 177, 16/23 on page 9 [ ]
%% Caret: " |G : N| `` r<Caret> is an odd p"
%% Comment: ""
%%
%⭣ ⭣ ⭣
$|G:N|=r$ is %%
%⭡ ⭡ ⭡ END of annotation 177
an
odd prime and $\theta\in \Irr(N)$ with %%
%% Annotation 175, 14/23 on page 9 [ ]
%% Caret: "Qpθq {\v D} Qpζpq<Caret> for some pr"
%% Comment: ""
%%
%⭣ ⭣ ⭣
$\QQ(\theta)\subseteq
\QQ(\zeta_p)$ for %%
%⭡ ⭡ ⭡ END of annotation 175
some prime $p\neq r$, then there exists an
irreducible constituent $\chi$ of $\theta^G$ such that
$\QQ(\chi)\subseteq \QQ(\zeta_p)$. This follows from
\cite[Lemma~2.2]{GHSV21}, for instance.

We restate Theorem~\ref{thm:main} for the reader's convenience.

\begin{theorem}
Let $\pi$ be a set of odd primes and $G$ a finite group such that
$G$ has a nontrivial Hall $\pi$-subgroup. Then $G$ possesses a
nontrivial $\pi'$-degree irreducible character with values %%
%% Annotation 174, 13/23 on page 9 [ ]
%% Caret: "ues in Qpζpq<Caret> for some p "
%% Comment: ""
%%
%⭣ ⭣ ⭣
in
$\QQ(\zeta_p)$ for %%
%⭡ ⭡ ⭡ END of annotation 174
some $p\in\pi$.
\end{theorem}

\begin{proof}
We argue by induction on $|G|$. %%
%% Annotation 173, 12/23 on page 9 [ ]
%% Caret: "|. Let M {\" Y} G<Caret> be a normal"
%% Comment: ""
%%
%⭣ ⭣ ⭣
Let $M \triangleleft\, G$ be %%
%⭡ ⭡ ⭡ END of annotation 173
a
normal subgroup such that $G/M$ is simple.

First assume %%
%% Annotation 172, 11/23 on page 9 [ ]
%% Caret: "ume that G{M<Caret> is abelian "
%% Comment: ""
%%
%⭣ ⭣ ⭣
that $G/M$ is %%
%⭡ ⭡ ⭡ END of annotation 172
abelian and that $|G/M| \in \{2\} \cup
\pi$. Then the inflation to $G$ of any nontrivial linear character
of $G/M$ satisfies the required conditions. Now suppose %%
%% Annotation 171, 10/23 on page 9 [ ]
%% Caret: "t r :`` |G{M|<Caret> is an odd p"
%% Comment: ""
%%
%⭣ ⭣ ⭣
that $r :=
|G/M|$ is %%
%⭡ ⭡ ⭡ END of annotation 171
an odd prime not belonging to $\pi$. In particular, we
have $\pi(M) \supseteq \pi\cap\pi(G)$.

By the induction hypothesis and the fact that $M$ also has a Hall
$\pi$-subgroup, there exists $\psi \in \Irr_{\pi'}(M)$ and %%
%% Annotation 170, 9/23 on page 9 [ ]
%% Caret: "d some p P π<Caret> such that Q"
%% Comment: ""
%%
%⭣ ⭣ ⭣
some $p
\in \pi$ such %%
%⭡ ⭡ ⭡ END of annotation 170
that $\QQ(\psi) \subseteq \QQ(\zeta_p)$. Let $\chi \in
\Irr(G)$ lie over $\psi$. %%
%% Annotation 166, 5/23 on page 9 [ ]
%% Caret: "`` {\v r}r  i``1 ψi<Caret> is the sum "
%% Comment: ""
%%
%% Annotation 167, 6/23 on page 9 [ ]
%% Caret: "er χM `` ψ or<Caret> χM `` {\v r}r  i``"
%% Comment: ""
%%
%% Annotation 168, 7/23 on page 9 [ ]
%% Caret: "ither χM `` ψ<Caret> or χM `` {\v r}r "
%% Comment: ""
%%
%% Annotation 169, 8/23 on page 9 [ ]
%% Caret: ".19], either<Caret> χM `` ψ or χ"
%% Comment: ""
%%
%⭣ ⭣ ⭣
By~\cite[Corollary~6.19]{Isaacs1}, either
$\chi_M = \psi$ or $\chi_M = \sum_{i=1}^r \psi_i$ is %%
%⭡ ⭡ ⭡ END of annotations 166, 167, 168, 169
the sum of the
$G$-conjugates of $\psi$.

In the latter case, we may take $\chi = \psi^G$, with a note %%
%% Annotation 162, 1/23 on page 9 [ ]
%% Caret: "χp1q `` rψp1q<Caret> is a π1- nu"
%% Comment: ""
%%
%⭣ ⭣ ⭣
that
$\chi(1) = r \psi(1)$ is %%
%⭡ ⭡ ⭡ END of annotation 162
%%
%% Annotation 163, 2/23 on page 9 [ ]
%% Caret: "p1q is a π1-<Caret> number and "
%% Comment: ""
%%
%% Annotation 164, 3/23 on page 9 [ ]
%% Caret: "p1q is a π1-<Caret> number and "
%% Comment: ""
%%
%% Annotation 165, 4/23 on page 9 [ ]
%% Caret: "rψp1q is a π<Caret>1- number an"
%% Comment: ""
%%
%⭣ ⭣ ⭣
a $\pi'$-number and %%
%⭡ ⭡ ⭡ END of annotations 163, 164, 165
$\QQ(\chi) \subseteq
\QQ(\psi)\subseteq \QQ(\zeta_p)$. In the former case, every
irreducible character of $G$ lying over $\psi$ is an extension of
$\psi$, and hence has $\pi'$-degree. Moreover, as mentioned above,
one of these extensions has values in $\QQ(\zeta_p)$, as required.

Finally, suppose that $G/M$ is non-abelian. Set $\pi_1 := \pi(G/M)
\cap \pi$. (Note that $\pi_1$ could be empty.) Again, the group
$G/M$ has a Hall $\pi$-subgroup, which is also a Hall
$\pi_1$-subgroup. Theorem~\ref{thm:simple} then yields a character
$\chi \in \Irr_{\pi_1'}(G/M)$ such that $\QQ(\chi) \subseteq
\QQ(\zeta_p)$ for some $p \in \pi_1$. Since $\chi(1)$ divides
$|G/M|$, the character $\chi$ is also of $\pi'$-degree. The desired
character is obtained by inflating $\chi$ to $G$.
\end{proof}

We conclude with a remark that, in view of the above proof, the
conclusion of Theorem~\ref{thm:main} remains valid when $2 \in \pi$,
provided that the group $G$ is $\pi$-separable.

%%%%%%%%%%%%%%%%%%%%%%
%%%%%%%%%%%%%%%%%%%%%%%

\begin{thebibliography}{ABCDEF}

\bibitem[B-S70]{Borel70}
A. Borel, R. W. Carter, C.\,W. Curtis, N. Iwahori, T.\,A. Springer,
and R. Steinberg, \emph{Seminar on algebraic groups and related
finite groups}, Lect. Notes in Math. 131, Springer-Verlag, Berlin,
1970.

\bibitem[Car85]{Carter85}
R.\,W. Carter, \emph{Finite groups of Lie type. Conjugacy classes
and complex characters}, Pure Appl. Math. (N. Y.) Wiley-Intersci.
Publ. John Wiley \& Sons, Inc., New York, 1985.

\bibitem[Atl]{Atl1}
  J.\,H. Conway, R.\,T. Curtis, S.\,P. Norton, R.\,A. Parker, and R.\,A. Wilson,
 {\it Atlas of finite groups}, Clarendon Press, Oxford, 1985.

%\bibitem[DM91]{Digne-Michel}
%  F. Digne and J. Michel, {\it Representations of finite
%  groups of Lie type}, London Mathematical Society Student Texts \textbf{95}, $2020$.

\bibitem[FT63]{Thompson63}
W. Feit and J. Thompson, Solvability of groups of odd order,
\emph{Pacific J. Math.} \textbf{13} (1963), 775--1029.

%\bibitem[GAP]{GAP}
%{The GAP Group}, \emph{GAP Groups, Algorithms, and Programming,
%Version 4.11.0}, 2020. \url{http://www.gap-system.org}




\bibitem[GM20]{GeckMalle}
M. Geck and G. Malle, {\it The character theory of finite groups of
Lie type: a guided tour}, Cambridge studies in advanced mathematics
\textbf{187}, Cambridge University Press, Cambridge, 2020.

\bibitem[GHSV21]{GHSV21}
%%
%% Annotation 194, 10/10 on page 10 [ ]
%% Caret: "] E. Giannel<Caret>li, N. N. Hu"
%% Comment: ""
%%
%⭣ ⭣ ⭣
E. Giannelli, N.%%
%⭡ ⭡ ⭡ END of annotation 194
\,N. Hung, %%
%% Annotation 185, 1/10 on page 10 [ ]
%% Caret: " A. Schaeffe<Caret>r Fry, and C"
%% Comment: ""
%%
%% Annotation 186, 2/10 on page 10 [ ]
%% Caret: "A. A. Schaef<Caret>fer Fry, and"
%% Comment: ""
%%
%% Annotation 187, 3/10 on page 10 [ ]
%% Caret: "A. A. Schaef<Caret>fer Fry, and"
%% Comment: ""
%%
%% Annotation 188, 4/10 on page 10 [ ]
%% Caret: "A. A. Schaef<Caret>fer Fry, and"
%% Comment: ""
%%
%% Annotation 189, 5/10 on page 10 [ ]
%% Caret: " A. A. Schae<Caret>ffer Fry, an"
%% Comment: ""
%%
%% Annotation 190, 6/10 on page 10 [ ]
%% Caret: ", A. A. Scha<Caret>effer Fry, a"
%% Comment: ""
%%
%% Annotation 191, 7/10 on page 10 [ ]
%% Caret: "g, A. A. Sch<Caret>aeffer Fry, "
%% Comment: ""
%%
%% Annotation 192, 8/10 on page 10 [ ]
%% Caret: "ng, A. A. Sc<Caret>haeffer Fry,"
%% Comment: ""
%%
%⭣ ⭣ ⭣
A.\,A. Schaeffer Fry, %%
%⭡ ⭡ ⭡ END of annotations 185, 186, 187, 188, 189, 190, 191, 192
and %%
%% Annotation 193, 9/10 on page 10 [ ]
%% Caret: ", and C. Val<Caret>lejo, Charac"
%% Comment: ""
%%
%⭣ ⭣ ⭣
C. Vallejo,
Characters %%
%⭡ ⭡ ⭡ END of annotation 193
of $\pi'$-degree and small cyclotomic fields, \emph{Ann.
Mat. Pura Appl.} \textbf{200} (2021), 1055--1073.

\bibitem[GLS98]{GLS98}
D. Gorenstein, R. Lyons, and R. Solomon, \emph{The classification of
the finite simple groups}, Number 3; Part I, Chapter A: Almost
simple $\mathcal{K}$-groups, Mathematical Surveys and Monographs,
40.3, American Mathematical Society, Providence, RI, 1998.

\bibitem[Gro86]{G86}
F. Gross, On a conjecture of Philip Hall, \emph{Proc. Lond. Math.
Soc. (3)} \textbf{52} (1986), no. 3, 464--494.

\bibitem[Hal28]{Hall1}
{P. Hall}, A note on soluble groups,
\emph{J. London. Math. Soc.} \textbf{3} (1928), 98--105.

\bibitem[Hal56]{Hall2}
{P. Hall}, Theorems like Sylow's,
\emph{Proc. London Math. Soc.} \textbf{6} (1956), 286--304.



\bibitem[Hun24]{Hung22}
{N.\,N. Hung}, The continuity of $p$-rationality and a lower bound
for $p'$-degree irreducible characters of finite groups,
\emph{Trans. Amer. Math. Soc.} \textbf{377} (2024), 323--344.

\bibitem[Isa06]{Isaacs1}
I.\,M. Isaacs, {\it Character theory of finite groups}, AMS Chelsea
Publishing, Providence, Rhode Island, 2006.

\bibitem[Mal07]{M07}
G. Malle, Height 0 characters of finite groups of Lie type,
\emph{Represent. Theory} \textbf{11} (2007), 192--220.

\bibitem[MT11]{malletesterman}
G.~Malle and D.~Testerman, \emph{Linear algebraic groups and finite
 groups of Lie type}, Cambridge studies in advanced mathematics
\textbf{133}, Cambridge University Press, Cambridge, 2011.


\bibitem[NT06]{Navarro-Tiep1}
G. Navarro and P.\,H. Tiep, Characters of $p'$-degree with
cyclotomic field of values, \emph{Proc. Amer. Math. Soc.}
\textbf{134} (2006), 2833--2837.

\bibitem[NT08]{Navarro-Tiep2}
 G. Navarro and P.\,H. Tiep, Rational irreducible
 characters and rational conjugacy classes in finite groups,
 \emph{Trans. Amer. Math. Soc.}
 \textbf{360} (2008), 2443--2465.

  \bibitem[Ras77]{Rasala}
R. Rasala, On the minimal degrees of characters of $\Sy_n$, \emph{J.
Algebra} \textbf{45} (1977), 132--181.

\bibitem[RV10]{RV10}
D.\,O. Revin and E.\,P. Vdovin, On the number of classes of
conjugate Hall subgroups in finite simple groups, \emph{J. Algebra}
\textbf{324} (2010), no. 12, 3614--3652.



\bibitem[Tho66]{Thompson}
J. Thompson, Hall subgroups of the symmetric groups,
\emph{J. Combinatorial Theory} \textbf{1} (1966), 271--279.

\end{thebibliography}
\end{document}
